% ============================================================
% Paper 1 patch snippets (β₁ stripped; graph language)
% These are drop-in replacements / insertions for main.tex.
% ============================================================

% ----------------------------
% Abstract insertion (replace the existing "cycle" formalization paragraph)
% ----------------------------
% Suggested replacement paragraph:
%
% We formalize this cycle as navigation in a model space $\mathcal{M}$ whose
% elements are candidate belief-dependency structures.  DLN stages are then
% characterized by two formal objects: the belief-dependency graph $G$
% (determining Level~1 inference cost) and a revision graph
% $\mathcal{R}=(\mathcal{M},\mathcal{T})$ over $\mathcal{M}$ (determining Level~2
% meta-cognitive capacity).  Linear cognition is a fixed point in $\mathcal{M}$
% with $\mathcal{T}=\emptyset$; Network cognition revises by traversing edges of
% $\mathcal{R}$, including a return transition that enables bounded recovery
% after model failure.

% ----------------------------
% Section 2.1 bridge sentence (replace the sentence that mentions "structural feature")
% ----------------------------
% Suggested replacement:
%
% Section~\ref{sec:model_space} formalizes this distinction by defining a model
% space $\mathcal{M}$ whose elements are candidate belief structures and a
% revision graph $\mathcal{R}$ whose edges are transitions between them; the
% cycle that defines Network cognition is a structural feature of $\mathcal{R}$,
% not of any individual belief graph $G\in\mathcal{M}$.

% ----------------------------
% Figure 1, Panel (b) label + caption edits
% ----------------------------
% In the Panel (b) TikZ, replace any line that annotates $a graph-structure annotation$
% with:
%
% \node[draw=none, below=0.7cm of gF3, anchor=west] {$\mathcal{R}$ contains a cycle};
%
% Caption replacement for Panel (b):
%
% (b) Revision graphs over model space (Level~2). Dot: empty model space.
% Linear: fixed point in model space (no revision edges). Network (expand-only):
% directed path $G_F \rightarrow G_{\mathrm{tab}}$. Network (full cycle): two
% directed transitions (expand and contract) forming a cycle in $\mathcal{R}$,
% enabling bounded recovery (Proposition~\ref{prop:crossover}(iii)).

% ----------------------------
% Discussion, Section 7.1 (replace deployment-layer paragraph + add meta-cognition paragraph)
% ----------------------------
% Replacement paragraph:
%
% The revision graph $\mathcal{R}$ introduced in Section~\ref{sec:model_space}
% provides a formal model of this deployment layer: the structural learning
% cycle implements Level~2 computation (navigation of $\mathcal{R}$), and the
% crossover condition $K^* = F + c_{\mathrm{meta}}/c_{\mathrm{param}}$ determines
% when investing in Level~2 is cost-effective.  This connects to metacognitive
% monitoring \cite{Flavell1979} by formalizing what is monitored (model adequacy,
% via predictive tests) and what actions are available in response (transitions
% in $\mathcal{R}$).
%
% New paragraph:
%
% \paragraph{Meta-cognition as a property of model space.}
% DLN stages differ not only in the graph-structure of the belief graph $G$ (Level~1)
% but in whether the agent can navigate a non-trivial model space (Level~2).
% Linear cognition is a fixed point in $\mathcal{M}$.  Network cognition
% traverses $\mathcal{M}$ via revision transitions, including a return transition
% whose functional consequence is bounded cost recovery after model failure
% (Proposition~\ref{prop:crossover}(iii)).  This provides a foundation for
% instantiating DLN in other domains: any system where an observer, agent, or
% learner can be characterized by (i)~a current representational structure and
% (ii)~a set of available transitions between structures can be classified by
% the structure of its revision graph.

% ----------------------------
% Conclusion replacement (replace the existing concluding paragraph that invokes graph-structure)
% ----------------------------
% Suggested replacement:
%
% Two formal objects carry the DLN distinction: the belief-dependency graph $G$,
% whose structure determines Level~1 inference cost, and the revision graph
% $\mathcal{R}$ over model space, whose structure determines Level~2
% meta-cognitive capacity.  Linear cognition is a fixed point in model space;
% Network cognition navigates $\mathcal{R}$ via hypothesis, test, and revision
% transitions.  Proposition~\ref{prop:crossover} derives the crossover
% $K^* = F + c_{\mathrm{meta}}/c_{\mathrm{param}}$ from the interaction between
% Level~1 compression and Level~2 overhead, and shows that including a return
% transition in $\mathcal{R}$ provides bounded recovery after model failure.
% This two-object characterization provides a foundation for instantiating DLN
% in domains beyond decision-making, including settings where representational
% constraints determine what information an agent can extract from its
% environment.
